\section{従来歩行制御手法の限界}

軟弱地面における歩行では，従来の剛体床を前提とした歩行制御手法の仮定が部分的に破綻する．多くの手法では，床変形による影響を外乱として扱い，フィードバック制御により補償することを想定している．しかし，床沈み込みは接地と同時に発生し，かつ継続的に変化する現象であるため，単純な外乱補償のみでは十分に対応できない．

また，CoM の水平方向運動の制御のみに着目した手法では，CoM 高さ変動に起因する ZMP 応答の変化を直接的に抑制することができない．その結果，安定化制御が過度に作用し，関節トルクの増大や不自然な運動を招く可能性がある．